\chapter[Operatory monotoniczne i p-laplasjan]{Operatory monotoniczne, twierdzenie Minty'ego--Browdera i nieliniowe równania eliptyczne}
\label{chap:pde26}

\section{Dlaczego metoda Laxa--Milgrama przestaje wystarczać?}
W rozdziale~\ref{chap:pde25} oraz we wcześniejszych zagadnieniach
eliptycznych występowały albo operatory liniowe, albo nieliniowości
niższego rzędu. Teraz najwyższa pochodna występuje nieliniowo:
\begin{equation}\label{eq:pde26-main}
 -\diver\bigl(|\nabla u|^{p-2}\nabla u\bigr)=f,\qquad p>1.
\end{equation}
Dla $p\ne2$ nie istnieje forma dwuliniowa pozwalająca zastosować
wprost twierdzenie Laxa--Milgrama. Zastępujemy dwuliniowość
\emph{monotonicznością}, która w teorii nieliniowej odgrywa rolę
dodatniości formy kwadratowej.

W całym rozdziale $\Omega\subset\R^d$ oznacza ograniczony obszar
Lipschitzowski, o ile przy konkretnym wyniku nie zaznaczono inaczej.
Ustalamy $1<p<\infty$, $p'=p/(p-1)$ oraz rzeczywistą
przestrzeń $V=W_0^{1,p}(\Omega)$ wyposażoną w normę
\[
 \|u\|_V=\left(\int_\Omega|\nabla u(x)|^p\,
                         \mathrm d\lambda_d(x)\right)^{1/p},
\]
gdzie moduł gradientu jest normą euklidesową.
Nierówność Poincarégo zapewnia jej równoważność ze zwykłą
normą Sobolewa: przedłużamy funkcje testowe przez zero
do kostki $Q$ zawierającej $\overline\Omega$ i stosujemy
Wniosek~\ref{cor:pde9-Poincare-zero}, po czym przechodzimy
do granicy w $W^{1,p}$. Odwzorowanie
$u\mapsto(u,\partial_1u,\ldots,\partial_du)$ identyfikuje
$W^{1,p}$ z domkniętą podprzestrzenią skończonego
iloczynu przestrzeni $L^p$ (domkniętość wynika
z definicji słabych pochodnych). Ponieważ $1<p<\infty$,
przestrzeń ta jest separowalna i refleksywna; te własności
dziedziczy domknięta podprzestrzeń $V$. Jej przestrzeń
sprzężoną oznaczamy przez $V^*=W^{-1,p'}(\Omega)$,
zawsze z normą funkcjonału względem normy $\|\cdot\|_V$.
Dla innej równoważnej normy Sobolewa w oszacowaniach
liczbowych może pojawić się stała równoważności.
Przez $\langle f,v\rangle$ oznaczamy wartość funkcjonału
$f\in V^*$ na $v\in V$; w przypadku $f\in L^{p'}(\Omega)$
jest to całka $\int_\Omega f(x)v(x)\,\mathrm d\lambda_d(x)$.
Dla funkcji $v\in V$ oznaczamy $v_+=\max\{v,0\}$ i
$v_-=\max\{-v,0\}$; obie należą do $V$, a
$\nabla v_+=\chi_{\{v>0\}}\nabla v$ prawie wszędzie
(Wniosek~\ref{cor:pde9-truncation}). Aby uzasadnić
zerowy warunek brzegowy, wybieramy $v_n\in C_c^\infty(\Omega)$
zbieżne do $v$ w $W^{1,p}$. Funkcje $(v_n)_+$
mają zwarty nośnik w $\Omega$ i należą do $W^{1,p}$,
więc po przedłużeniu przez zero można je przybliżyć
splotami o nośnikach nadal zawartych w $\Omega$.
Tym samym należą do $V$. Są ograniczone w $V$
dzięki wzorowi na gradient, a w $L^p$ zbiegają do $v_+$
z nierówności $|a_+-b_+|\le|a-b|$.
Z refleksywności $V$ wybieramy podciąg słabo zbieżny
w $V$ do pewnego $w\in V$. Ciągłe zanurzenie
$V\hookrightarrow L^p$ daje jego słabą zbieżność
do $w$ również w $L^p$. Jednoznaczność granicy
słabej wymusza $w=v_+$ jako klasy funkcji.
Tak samo postępujemy dla $v_-=(-v)_+$.

\begin{definicja}[Monotoniczność i hemiciągłość]\label{def:pde26-monotone}
Niech $X$ będzie rzeczywistą przestrzenią Banacha, a
$A:X\to X^*$ odwzorowaniem określonym na całym $X$.
Mówimy, że $A$ jest
\emph{monotoniczne}, jeżeli
$\langle Au-Av,u-v\rangle\ge0$ dla wszystkich $u,v\in X$,
zaś \emph{ściśle monotoniczne}, jeżeli nierówność jest ostra
przy $u\ne v$. Operator jest \emph{hemiciągły}, jeśli dla
wszystkich $u,v,w\in X$ odwzorowanie
$t\mapsto\langle A(u+tv),w\rangle$ jest ciągłe na $\R$.
Operator $A$ jest \emph{ograniczony na zbiorach ograniczonych},
jeśli dla każdego $R>0$ zachodzi
$\sup_{\|u\|_X\le R}\|Au\|_{X^*}<\infty$.
Nazywamy go \emph{koercywnym}, jeśli
$\langle Au,u\rangle/\|u\|_X\to+\infty$
przy $\|u\|_X\to\infty$.
\end{definicja}

\begin{uwaga}[Warianty słowa koercywność]
Spotyka się również warunek silniejszy
$\langle Au,u\rangle\ge c\|u\|_X^p-C$ dla
$c>0$, $C\ge0$ i $p>1$.
Jest on wystarczający do naszej definicji, ale nie jest z nią
równoważny: dla $X=\R$ i $A(s)=\operatorname{sgn}(s)
\log(1+|s|)$ iloraz koercywności wynosi $\log(1+|s|)$,
natomiast $\langle A(s),s\rangle=|s|\log(1+|s|)$
rośnie wolniej niż $|s|^p$ dla każdego $p>1$.
Koercywność nie oznacza monotoniczności:
$A(s)=s^3-s$ jest koercywny, lecz
$(A(1/2)-A(0))(1/2-0)<0$.
Z kolei operator zerowy jest monotoniczny, ale nie
jest koercywny na $\R$. Obie własności są zatem
niezależnymi hipotezami twierdzenia Minty'ego--Browdera.
\end{uwaga}

\section{Nierówności monotoniczności dla wektorów}
\begin{lemat}[Monotoniczność strumienia potęgowego]\label{lem:pde26-vector}
Dla $p>1$ odwzorowanie $a_p:\R^d\to\R^d$,
$a_p(\xi)=|\xi|^{p-2}\xi$ dla $\xi\ne0$ i $a_p(0)=0$,
jest ciągłe i ściśle monotoniczne. Co więcej,
dla wszystkich $\xi,\eta\in\R^d$ zachodzą:
\begin{align}
 (a_p(\xi)-a_p(\eta))\cdot(\xi-\eta)
 &\ge \frac{2^{2-p}}{p-1}|\xi-\eta|^p &&(p\ge2),\label{eq:pde26-strong}\\
\intertext{a dla $1<p<2$ i $\xi\ne\eta$}
 (a_p(\xi)-a_p(\eta))\cdot(\xi-\eta)
 &\ge(p-1)(|\xi|+|\eta|)^{p-2}|\xi-\eta|^2.\label{eq:pde26-subquad}
\end{align}
W obu przypadkach prawa strona jest dodatnia przy $\xi\ne\eta$.
\end{lemat}
\begin{proof}
Funkcja $\Phi(\xi)=|\xi|^p/p$ jest ściśle wypukła,
a jej gradientem jest $a_p$. Wypukłość daje monotoniczność;
wyjaśniamy dodatkowo jej ilościowy charakter.
Dla $\zeta\ne0$ macierz pochodnej $a_p$ wynosi
$Da_p(\zeta)=|\zeta|^{p-2}I+(p-2)|\zeta|^{p-4}\zeta\otimes\zeta$.
W kierunku $\zeta$ macierz ma wartość własną
$(p-1)|\zeta|^{p-2}$, a w kierunkach prostopadłych
$|\zeta|^{p-2}$, zatem jest ograniczona od dołu
w sensie form kwadratowych przez
$\min\{1,p-1\}|\zeta|^{p-2}I$ (także dla $d=1$).
Dla $\xi=\eta$ nierówność monotoniczności jest trywialna.
Niech $h=\xi-\eta\ne0$. Jeżeli odcinek
$\eta+th$ zawiera zero dla $t=t_*$, to
$|\eta+th|=|h||t-t_*|$. Pochodna funkcji
$t\mapsto a_p(\eta+th)$ jest wtedy ograniczona
w normie przez stałą razy $|t-t_*|^{p-2}$,
co jest całkowalne dla $p>1$.
Stosując twierdzenie Newtona--Leibniza oddzielnie
po obu stronach $t_*$ i przechodząc do granicy
dzięki ciągłości $a_p$ w zerze, uzyskujemy ten sam
wzór całkowy co dla odcinka omijającego zero.
Otrzymujemy
\[
 (a_p(\xi)-a_p(\eta))\cdot(\xi-\eta)
 \ge \min\{1,p-1\}|\xi-\eta|^2
 \int_0^1|\eta+t(\xi-\eta)|^{p-2}\,\mathrm d\lambda_1(t).
\]
Jeśli $1<p<2$, na całym odcinku
$|\eta+t(\xi-\eta)|\le|\xi|+|\eta|$, a wykładnik
$p-2$ jest ujemny. Otrzymujemy~\eqref{eq:pde26-subquad}.
Dla $p=2$ mamy dokładnie $a_2(\xi)=\xi$.
Dla $p>2$ korzystamy z elementarnej nierówności
\[
 \int_0^1|\eta+t(\xi-\eta)|^{p-2}\,\mathrm d\lambda_1(t)
 \ge \frac{2^{2-p}}{p-1}|\xi-\eta|^{p-2}.
\]
Dla przejrzystości można ją uzasadnić po rozłożeniu
$\eta$ na składową równoległą i prostopadłą do
$h=\xi-\eta$: usunięcie składowej prostopadłej
nie zwiększa lewej strony, a funkcja
$b\mapsto\int_0^1|b+t|^{p-2}\,\mathrm d\lambda_1(t)$
ma pochodną $|b+1|^{p-2}-|b|^{p-2}$,
na mocy twierdzenia Newtona--Leibniza.
Jest ona ujemna dla $b<-1/2$ i dodatnia dla
$b>-1/2$, ponieważ $p-2>0$. Minimum osiągamy
więc w $b=-1/2$ i wynosi ono
$2\int_0^{1/2}t^{p-2}\,\mathrm d\lambda_1(t)
=2^{2-p}/(p-1)$.
To prowadzi do stałej $2^{2-p}/(p-1)$ w
\eqref{eq:pde26-strong}. Nie jest to najlepsza możliwa
stała, ale wystarcza do dalszych zastosowań.
Ścisłość wynika z dodatniości uzyskanych prawych stron.
\end{proof}

\begin{przyklad}[Dlaczego $p=2$ jest szczególny]
Dla $p=2$ mamy $a_2(\xi)=\xi$ i
$(a_2(\xi)-a_2(\eta))\cdot(\xi-\eta)=|\xi-\eta|^2$.
Dla $p>2$ różniczka $Da_p(0)$ jest zerowa, zatem
jednostajna eliptyczność linearyzacji zanika przy
$\nabla u=0$. Dla $1<p<2$ pochodna $a_p$ ma
osobliwość przy zerze. Żadne z tych zjawisk nie
niszczy jednak monotoniczności nieliniowego operatora.
\end{przyklad}

\section{Twierdzenie Minty'ego--Browdera z dowodem}
\begin{lemat}[Skończeniewymiarowa zasada istnienia]
\label{lem:pde26-brouwer}
Jeśli $F:\R^n\to\R^n$ jest ciągłe, a dla pewnego
$R>0$ zachodzi $F(x)\cdot x>0$ na sferze $|x|=R$,
to istnieje $x\in\R^n$ takie, że $|x|<R$ i $F(x)=0$.
\end{lemat}
\begin{proof}
Załóżmy przeciwnie, że $F$ nie zeruje się na
$\overline B(0,R)$. Definiujemy na tej kuli
odwzorowanie ciągłe
\[
 T(x)=-\frac{RF(x)}{|F(x)|}.
\]
Jego wartości należą do sfery o promieniu $R$,
a więc do rozważanej kuli. Twierdzenie Brouwera
o punkcie stałym daje $x=-RF(x)/|F(x)|$.
Wtedy $|x|=R$ oraz $F(x)\cdot x=-R|F(x)|<0$,
co jest sprzeczne z hipotezą. Wniosek jest prawdziwy.
\end{proof}

\begin{twierdzenie}[Twierdzenie Minty'ego--Browdera; wersja z pełną dziedziną]
\label{thm:pde26-minty}
Niech $X$ będzie rzeczywistą, separowalną, refleksywną
przestrzenią Banacha. Załóżmy, że $A:X\to X^*$ jest
monotoniczny, hemiciągły, ograniczony na zbiorach ograniczonych
i koercywny w sensie definicji~\ref{def:pde26-monotone}.
Wówczas dla każdego $f\in X^*$ istnieje $u\in X$ takie, że
$Au=f$. Jeśli ponadto $A$ jest ściśle monotoniczny,
to rozwiązanie jest jedyne.
\end{twierdzenie}
\begin{proof}
\textbf{Krok 1: przybliżenia Galerkina.}
Jeżeli $X=\{0\}$, teza jest natychmiastowa.
W przypadku nieskończenie wymiarowym wybieramy
przeliczalny gęsty podzbiór $X$ i usuwamy kolejno
elementy należące do powłoki liniowej poprzednich.
Otrzymujemy liniowo niezależny ciąg $(e_j)$
o gęstej powłoce liniowej. Kładziemy
$X_n=\operatorname{span}\{e_1,\ldots,e_n\}$.
Dla skończenie wymiarowego $X$ wybieramy bazę
i od chwili osiągnięcia jej wymiaru przyjmujemy
$X_n=X$; w poniższym układzie używamy bazy $X_n$.
Szukamy $u_n\in X_n$ spełniającego
\begin{equation}\label{eq:pde26-galerkin}
 \langle Au_n,v\rangle=\langle f,v\rangle
 \quad\text{dla każdego }v\in X_n.
\end{equation}
Oznaczmy $m_n=\dim X_n$. W bazie
$(e_1,\ldots,e_{m_n})$ tworzymy układ
$F:\R^{m_n}\to\R^{m_n}$, którego współrzędne to
\[
 F_j(c)=\left\langle A\left(\sum_{k=1}^{m_n}c_ke_k\right)
                      -f,e_j\right\rangle.
\]
Równość $F(c)=0$ jest równoważna
~\eqref{eq:pde26-galerkin} dla $u_n=\sum_{k=1}^{m_n}c_ke_k$,
gdyż funkcjonały są liniowe względem argumentu testowego.
Hemiciągłość wraz z monotonicznością i lokalną
ograniczonością implikuje ciągłość $A$ od normy $X$
do słabej topologii $X^*$: istotnie, jeśli $z_m\to z$,
to każda słaba granica $Az_m$ (istniejąca na podciągu)
spełnia, po przejściu do granicy w nierówności
$\langle Az_m-Av,z_m-v\rangle\ge0$,
relację $\langle g-Av,z-v\rangle\ge0$ dla wszystkich $v$.
Dokładniej, rozkładamy iloczyn na
$\langle Az_m-Av,z-v\rangle+
\langle Az_m-Av,z_m-z\rangle$.
Pierwszy składnik zbiega do podanej wartości,
a moduł drugiego jest ograniczony przez
$(M+\|Av\|_{X^*})\|z_m-z\|_X\to0$.
Teraz,
podstawiając $v=z-tw$ i przechodząc z $t\downarrow0$
przez hemiciągłość, otrzymujemy
$\langle g-Az,w\rangle\ge0$. Podstawienie $-w$
daje równość, a więc $g=Az$.
Ponieważ $(Az_m)$ jest ograniczony i $X^*$ refleksywna,
każdy podciąg ma słabo zbieżny dalszy podciąg.
Gdyby dla pewnego $w\in X$ ciąg
$\langle Az_m,w\rangle$ nie zbiegał do
$\langle Az,w\rangle$, można by wybrać podciąg
oddalony od tej wartości o dodatnią stałą,
co przeczyłoby właśnie wykazanej identyfikacji.
Stąd wszystkie współrzędne $F_j$ są ciągłe,
a więc $F$ jest ciągłe.

Koercywność i nierówność
$\langle Au-f,u\rangle\ge\langle Au,u\rangle
-\|f\|_{X^*}\|u\|_X$ pokazują, że
$\langle Au-f,u\rangle>0$ dla wszystkich
$\|u\|_X\ge R$ przy dostatecznie dużym $R$,
niezależnym od $n$.
Dla każdego $n$ równoważność norm na $X_n$ daje
stałą $b_n>0$ taką, że
$\|\sum_{j=1}^{m_n} c_je_j\|_X\ge b_n|c|$.
Wybieramy $r_n=R/b_n$, więc norma elementu $X_n$
jest co najmniej $R$ dla $|c|=r_n$.
Wówczas $F(c)\cdot c=\langle Au-f,u\rangle>0$.
Lemat~\ref{lem:pde26-brouwer} daje rozwiązanie
\eqref{eq:pde26-galerkin}.

\textbf{Krok 2: oszacowania i granice słabe.}
Testując~\eqref{eq:pde26-galerkin} przez $u_n$, mamy
$\langle Au_n,u_n\rangle=\langle f,u_n\rangle$.
Koercywność wymusza jednostajne ograniczenie
$\|u_n\|_X\le R$. Lokalna ograniczoność $A$ daje
$\|Au_n\|_{X^*}\le M_R$. Refleksywność $X$ oraz $X^*$
oraz twierdzenie Eberleina--Šmuliana pozwalają
wybrać rosnący ciąg indeksów $(n_j)$, dla którego
$u_{n_j}\rightharpoonup u$ w $X$ oraz
$Au_{n_j}\rightharpoonup g$ w $X^*$.
W dalszych granicach zapisujemy nadal $n$,
rozumiejąc przez niego indeks z tego podciągu,
a $X_n$ oznacza odpowiadającą mu pierwotną podprzestrzeń.
Dla dowolnego $j$ i $n\ge j$ mamy
$\langle Au_n,e_j\rangle=\langle f,e_j\rangle$.
Po przejściu do granicy i wykorzystaniu gęstości
otrzymujemy $g=f$.

\textbf{Krok 3: kluczowa nierówność graniczna.}
Wybieramy $v_n\in X_n$ takie, że $v_n\to u$ silnie
w $X$. Istnienie takiego ciągu wynika z gęstości
rosnących podprzestrzeni $X_n$: dla każdego $k$ wybieramy
$z_k\in X_{n_k}$ z $\|z_k-u\|_X<1/k$ i
$n_k\uparrow\infty$; kładziemy $v_n=z_k$ dla
$n_k\le n<n_{k+1}$. Dzięki równaniu Galerkina
$\langle Au_n,u_n-v_n\rangle=\langle f,u_n-v_n\rangle$.
Ponieważ $u_n\rightharpoonup u$ i $v_n\to u$,
mamy $\langle f,u_n-v_n\rangle\to0$.
Ponadto
$|\langle Au_n,v_n-u\rangle|\le M_R\|v_n-u\|_X\to0$.
Dodając obie równości, otrzymujemy
\begin{equation}\label{eq:pde26-limsup}
 \lim_{n\to\infty}\langle Au_n,u_n-u\rangle=0.
\end{equation}
Dla ustalonego $v\in X$ monotoniczność daje
$\langle Au_n-Av,u_n-v\rangle\ge0$.
Rozkładamy pierwszy człon jako
$\langle Au_n,u_n-u\rangle+\langle Au_n,u-v\rangle$
i korzystamy z~\eqref{eq:pde26-limsup}. W granicy
otrzymujemy $\langle f-Av,u-v\rangle\ge0$.

\textbf{Krok 4: identyfikacja operatora.}
Dla dowolnego $w\in X$ i $t>0$ podstawiamy $v=u-tw$.
Wynika stąd $\langle f-A(u-tw),w\rangle\ge0$.
Przy $t\downarrow0$ hemiciągłość daje
$\langle f-Au,w\rangle\ge0$. Stosując nierówność także
do $-w$, otrzymujemy $Au=f$.
Jeśli $Au=Av=f$, ścisła monotoniczność wymusza $u=v$.
\end{proof}

\begin{uwaga}[Dokładny zakres twierdzenia]
Nie wymagamy liniowości $A$ ani istnienia potencjału
wariacyjnego. Warunek ograniczoności na zbiorach
ograniczonych jest tu przyjęty jawnie, ponieważ używamy
go w dowodzie do wydzielenia słabo zbieżnego podciągu $Au_n$.
Twierdzenie nie obejmuje automatycznie operatorów
określonych jedynie na właściwej części $X$ ani równań,
w których zawodzi monotoniczność.
\end{uwaga}

\section{p-laplasjan jako operator na przestrzeni Sobolewa}
\begin{definicja}[Rozwiązanie słabe $p$-laplasjanu]\label{def:pde26-weak}
Niech $f\in V^*$. Funkcję $u\in V=W_0^{1,p}(\Omega)$
nazywamy słabym rozwiązaniem zagadnienia
\eqref{eq:pde26-main} z zerowym warunkiem Dirichleta,
jeśli dla każdego $v\in V$ zachodzi
\begin{equation}\label{eq:pde26-weak}
 \int_\Omega|\nabla u|^{p-2}\nabla u\cdot\nabla v
 \,\mathrm d\lambda_d(x)=\langle f,v\rangle.
\end{equation}
Operator $A_p:V\to V^*$ definiujemy przez lewą stronę
tego wzoru; zapis $-\Delta_pu=f$ oznacza równość
$A_pu=f$ w $V^*$.
\end{definicja}

\begin{twierdzenie}[Istnienie i jednoznaczność]\label{thm:pde26-plap}
Niech $1<p<\infty$ i niech $\Omega\subset\R^d$ będzie
ograniczonym obszarem Lipschitzowskim. Dla każdego
$f\in W^{-1,p'}(\Omega)$ istnieje dokładnie jedno
$u\in W_0^{1,p}(\Omega)$ spełniające
\eqref{eq:pde26-weak}. Ponadto
\begin{equation}\label{eq:pde26-apriori}
 \|\nabla u\|_{L^p(\Omega)}^{p-1}
 \le\|f\|_{W^{-1,p'}(\Omega)}.
\end{equation}
\end{twierdzenie}
\begin{proof}
Sprawdzamy każdą hipotezę twierdzenia
\ref{thm:pde26-minty} dla $A_p$.
Nierówność Höldera daje
$|\langle A_pu,v\rangle|\le
\|\nabla u\|_p^{p-1}\|\nabla v\|_p$, więc
$\|A_pu\|_{V^*}\le\|u\|_V^{p-1}$.
Dla $u\ne0$ test $v=u/\|u\|_V$ osiąga to ograniczenie:
$\langle A_pu,v\rangle=\|u\|_V^{p-1}$.
Dla $u=0$ obie strony są zerowe, zatem
$\|A_pu\|_{V^*}=\|u\|_V^{p-1}$ dla każdego $u\in V$.
Operator jest określony na całym $V$ i ograniczony
na zbiorach ograniczonych. Ciągłość strumienia $a_p$ oraz
oszacowanie
$|a_p(\nabla u+t\nabla w)|\le
C_p(|\nabla u|+|\nabla w|)^{p-1}$ dla $|t|\le1$
po pomnożeniu przez $|\nabla v|$ dają całkowalną
majorantę: $(p-1)p'=p$, więc stosujemy nierówność
Höldera. Twierdzenie o zbieżności zmajoryzowanej
zapewnia ciągłość $\langle A_p(u+tw),v\rangle$
w $t=0$, a po zastąpieniu $u$ przez $u+t_0w$
również w dowolnym $t_0\in\R$.
Monotoniczność wynika z Lematu~\ref{lem:pde26-vector}.
Jeżeli $u\ne v$ w $V$, to z nierówności Poincarégo
gradienty różnią się na zbiorze dodatniej miary,
na którym całkowany iloczyn jest ściśle dodatni.
Stąd operator jest ściśle monotoniczny.
Koercywność jest dokładna:
$\langle A_pu,u\rangle=\|\nabla u\|_p^p$,
a więc iloraz przez $\|u\|_V$ dąży do nieskończoności.
Twierdzenie Minty'ego--Browdera zapewnia jedyne rozwiązanie.
Testując jego równanie przez $u$, otrzymujemy
$\|u\|_V^p=\langle f,u\rangle\le
\|f\|_{V^*}\|u\|_V$. Dzielenie przez $\|u\|_V$
(dla $u=0$ teza jest oczywista) kończy dowód.
\end{proof}

\section{Metoda bezpośrednia rachunku wariacyjnego}
\begin{twierdzenie}[Charakteryzacja minimizacyjna]\label{thm:pde26-minimizer}
Niech $p>1$, $\Omega$ będzie ograniczonym obszarem
Lipschitzowskim, a $f\in V^*$. Funkcjonał
\[
 J_f(v)=\frac1p\int_\Omega|\nabla v|^p
 \,\mathrm d\lambda_d(x)-\langle f,v\rangle,\qquad v\in V,
\]
ma dokładnie jeden punkt minimum. Jest nim jedyne
rozwiązanie zagadnienia~\eqref{eq:pde26-weak}.
\end{twierdzenie}
\begin{proof}
Z definicji normy funkcjonału otrzymujemy
$J_f(v)\ge\|v\|_V^p/p-\|f\|_{V^*}\|v\|_V$.
Ponieważ $p>1$, prawa strona dąży do nieskończoności
przy $\|v\|_V\to\infty$ i jest ograniczona od dołu
jako ciągła funkcja nieujemnej zmiennej $\|v\|_V$.
Ciąg minimalizujący jest ograniczony w refleksywnej
przestrzeni $V$; wybrany podciąg słabo zbiega do $u$.
Ciągły operator liniowy $\nabla:V\to L^p(\Omega;\R^d)$
zachowuje zbieżność słabą. Norma w przestrzeni $L^p$
jest słabo półciągła z dołu (jako supremum modułów
ciągłych funkcjonałów liniowych o normie co najwyżej $1$),
a funkcjonał liniowy $v\mapsto\langle f,v\rangle$
jest słabo ciągły. Stąd $J_f(u)\le\liminf J_f(v_n)$,
czyli $u$ jest minimizatorem. Ścisła wypukłość
$z\mapsto|z|^p$ daje ścisłą wypukłość $J_f$ na $V$:
jeżeli $u\ne v$ w $V$, nierówność Poincarégo
zapewnia, że gradienty różnią się na zbiorze dodatniej
miary. Dla $0<\theta<1$ całkowanie punktowej nierówności
ścisłej wypukłości daje zatem
$J_f((1-\theta)u+\theta v)<(1-\theta)J_f(u)+\theta J_f(v)$.
Stąd minimum jest jedyne.

Dla dowolnych $u,v\in V$ i $|t|\le1$ iloraz różnicowy
funkcji $|\nabla u+t\nabla v|^p/p$ jest ograniczony
w module przez $(|\nabla u|+|\nabla v|)^{p-1}|\nabla v|$,
co wynika z twierdzenia o wartości średniej.
Majoranta jest całkowalna dzięki nierówności Höldera.
Z twierdzenia o zbieżności zmajoryzowanej wynika
\[
 \left.\frac{\mathrm d}{\mathrm dt}J_f(u+tv)\right|_{t=0}
 =\langle A_pu-f,v\rangle.
\]
Poniżej uzasadnimy również, że $J_f$ jest klasy $C^1$
w sensie różniczkowania Frécheta.
Dla minimizatora pochodna ta jest zerowa w każdym
kierunku, więc zachodzi~\eqref{eq:pde26-weak}.
Odwrotnie, jeżeli $A_pu=f$, to dla dowolnego $w\in V$
wypukła funkcja $t\mapsto J_f(u+t(w-u))$ spełnia
$J_f(w)-J_f(u)\ge\langle A_pu-f,w-u\rangle=0$.
Dowodzi to minimalności.
\end{proof}

\begin{uwaga}[Dlaczego przedstawiliśmy dwa dowody?]
Metoda minimizacji wykorzystuje szczególną strukturę
potencjałową: pochodna kierunkowa $J_f$ w $u$
jest dana przez funkcjonał $A_pu-f$. Twierdzenie Minty'ego--Browdera
działa również dla monotonicznych operatorów bez potencjału.
Przykładem jest strumień $a(x,\xi)$ monotoniczny w $\xi$,
który nie musi być gradientem funkcji wypukłej względem $\xi$.
\end{uwaga}


\begin{twierdzenie}[Strumienie typu Leraya--Lionsa]
\label{thm:pde26-leraylions}
Niech $\Omega\subset\R^d$ będzie ograniczonym obszarem
Lipschitzowskim, $1<p<\infty$ i $V=W_0^{1,p}(\Omega)$.
Niech $a:\Omega\times\R^d\to\R^d$ będzie funkcją
Carathéodory'ego: $x\mapsto a(x,\xi)$ jest mierzalna
przy ustalonym $\xi$, a $\xi\mapsto a(x,\xi)$ jest
ciągła dla prawie każdego $x$. Załóżmy, że istnieją
$C,\alpha>0$, $h\in L^{p'}(\Omega)$, $h\ge0$,
i $\beta\in L^1(\Omega)$, $\beta\ge0$, takie że
prawie wszędzie w $x$ i dla wszystkich $\xi,\eta\in\R^d$
zachodzą
\begin{align}
 |a(x,\xi)|&\le h(x)+C|\xi|^{p-1},\label{eq:pde26-ll-growth}\\
 a(x,\xi)\cdot\xi&\ge\alpha|\xi|^p-\beta(x),\label{eq:pde26-ll-coerc}\\
 (a(x,\xi)-a(x,\eta))\cdot(\xi-\eta)&\ge0.\label{eq:pde26-ll-mono}
\end{align}
Wówczas dla każdego $f\in V^*$ istnieje $u\in V$
spełniające
\begin{equation}\label{eq:pde26-ll-weak}
 \int_\Omega a(x,\nabla u(x))\cdot\nabla v(x)
 \,\mathrm d\lambda_d(x)=\langle f,v\rangle
 \qquad(v\in V).
\end{equation}
Jeżeli nierówność w~\eqref{eq:pde26-ll-mono} jest ostra
przy $\xi\ne\eta$ dla prawie każdego $x$, rozwiązanie
jest jednoznaczne.
\end{twierdzenie}
\begin{proof}
Definiujemy $A:V\to V^*$ lewą stroną
\eqref{eq:pde26-ll-weak}. Najpierw zauważmy, że dla
dowolnej mierzalnej funkcji $\zeta:\Omega\to\R^d$
złożenie $x\mapsto a(x,\zeta(x))$ jest mierzalne
w sensie Lebesgue'a: przybliżamy $\zeta$ punktowo
mierzalnymi funkcjami prostymi $\zeta_n$ o skończonym
zbiorze wartości. Funkcje $a(x,\zeta_n(x))$ są mierzalne
na mocy mierzalności przekrojów, a poza wspólnym
zbiorem wyjątkowym ciągłość względem $\xi$ daje
ich punktową zbieżność do $a(x,\zeta(x))$.
Stosujemy to do $\zeta=\nabla u$.
Zmiana reprezentanta gradientu na zbiorze miary zero
nie zmienia klasy otrzymanego strumienia.
Warunek wzrostu i nierówność Höldera zapewniają
$\|Au\|_{V^*}\le\|h\|_{p'}+C\|u\|_V^{p-1}$;
odwzorowanie jest więc dobrze określone i ograniczone
na zbiorach ograniczonych. Monotoniczność wynika
z całkowania~\eqref{eq:pde26-ll-mono}.
Dla $u,v,w\in V$ oraz $|t|\le1$ zachodzi prawie wszędzie
\[
 |a(x,\nabla u+t\nabla v)\cdot\nabla w|
 \le\bigl(h(x)+C(|\nabla u|+|\nabla v|)^{p-1}\bigr)
 |\nabla w|.
\]
Prawa strona jest całkowalna przez nierówność Höldera.
Ciągłość w zmiennej $\xi$ i twierdzenie o zbieżności
zmajoryzowanej dowodzą ciągłości w $t=0$.
Aby uzyskać ciągłość w dowolnym $t_0\in\R$,
zastępujemy $u$ przez $u+t_0v$.

W rzeczywistości $A:V\to V^*$ jest ciągły w normie.
Niech $u_n\to u$ w $V$. Z każdego podciągu wybieramy
dalszy podciąg $(u_{n_j})$ tak, by
$\|\nabla u_{n_j}-\nabla u\|_p\le 2^{-j}$.
Funkcja
\[
 H=|\nabla u|+\sum_{j=1}^\infty
                         |\nabla u_{n_j}-\nabla u|
\]
należy do $L^p$ na mocy nierówności Minkowskiego
dla sum częściowych i lematu Fatou.
W szczególności szereg jest skończony prawie wszędzie,
więc $\nabla u_{n_j}\to\nabla u$ prawie wszędzie.
Ponadto $|\nabla u_{n_j}|\le H$. Z warunku wzrostu
wynika oszacowanie
\[
 |a(x,\nabla u_{n_j})-a(x,\nabla u)|^{p'}
 \le (2h(x)+2CH(x)^{p-1})^{p'}.
\]
Prawa strona jest całkowalna, ponieważ $(p-1)p'=p$.
Ciągłość strumienia względem $\xi$ i twierdzenie
o zbieżności zmajoryzowanej dają zbieżność strumieni
w $L^{p'}(\Omega;\R^d)$ wzdłuż dalszego podciągu.
Z nierówności Höldera wynika
$\|Au_{n_j}-Au\|_{V^*}\to0$.
Gdyby cały ciąg nie zbiegał do $Au$, można by
wybrać podciąg oddalony w normie od $Au$ o dodatnią
stałą, wbrew właśnie udowodnionej własności.
Z~\eqref{eq:pde26-ll-coerc} otrzymujemy
$\langle Au,u\rangle\ge\alpha\|u\|_V^p-\|\beta\|_1$,
a zatem $A$ jest koercywny.
Wszystkie hipotezy twierdzenia~\ref{thm:pde26-minty}
są spełnione, więc istnieje rozwiązanie.
Dla dwóch rozwiązań $u_1,u_2$ testujemy różnicę
równań przez $u_1-u_2$. Jeśli monotoniczność jest
ostra, nieujemna całkowana funkcja może mieć całkę
zero tylko gdy $\nabla u_1=\nabla u_2$ prawie wszędzie.
Nierówność Poincarégo daje $u_1=u_2$.
\end{proof}

\begin{uwaga}[Różniczka funkcjonału energii]
Strumień $a(x,\xi)=a_p(\xi)$ spełnia hipotezy
Twierdzenia~\ref{thm:pde26-leraylions} z $h=\beta=0$,
$C=\alpha=1$. Z dowodu wynika zatem ciągłość
$A_p:V\to V^*$ w normie.
Dla $w\in V$ stosujemy twierdzenie Newtona--Leibniza
do funkcji $t\mapsto J_f(u+tw)$, której pochodną
obliczyliśmy w dowodzie Twierdzenia~\ref{thm:pde26-minimizer}.
Otrzymujemy
\[
 \begin{aligned}
 &|J_f(u+w)-J_f(u)-\langle A_pu-f,w\rangle|\\
 &\quad\le\|w\|_V\sup_{0\le t\le1}
                         \|A_p(u+tw)-A_pu\|_{V^*}.
 \end{aligned}
\]
Supremum dąży do zera przy $\|w\|_V\to0$, ponieważ
$\|tw\|_V\le\|w\|_V$ i $A_p$ jest ciągły w $u$.
W konsekwencji $J_f$ jest różniczkowalny w sensie
Frécheta, $DJ_f(u)=A_pu-f$, a różniczka jest ciągła.
\end{uwaga}

\begin{przyklad}[Monotoniczny strumień bez potencjału]
\label{ex:pde26-skewflux}
Niech $d=2$, $p=2$ i
$J=\begin{psmallmatrix}0&-1\\1&0\end{psmallmatrix}$.
Dla niekonstantnej funkcji $\kappa\in W^{1,\infty}(\Omega)$
kładziemy $a(x,\xi)=(I+\kappa(x)J)\xi$.
Wówczas $(a(x,\xi)-a(x,\eta))\cdot(\xi-\eta)
=|\xi-\eta|^2$ oraz $a(x,\xi)\cdot\xi=|\xi|^2$.
Twierdzenie~\ref{thm:pde26-leraylions} zapewnia jednoznaczne
rozwiązanie dla każdego $f\in H^{-1}(\Omega)$.
Na zbiorze $\{\kappa\ne0\}$ macierz $I+\kappa J$
nie jest symetryczna, zatem strumień $a(x,\cdot)$
nie jest tam gradientem funkcji skalarnej względem
$\xi$: gdyby był, ta funkcja byłaby klasy $C^2$
i jej hesjan byłby symetryczny.
Dla gładkich funkcji testowych rachunek daje
\[
 -\diver(\kappa J\nabla u)
 =-\nabla\kappa\cdot J\nabla u,
\]
gdyż $\diver(J\nabla u)=-\partial_1\partial_2 u
+\partial_2\partial_1 u=0$.
Dla przykładu $\kappa(x)=x_1$ otrzymujemy dodatkowy
człon $\partial_2 u$. Rachunek obowiązuje także
dystrybucyjnie dla $u\in H_0^1(\Omega)$, przez
aproksymację funkcjami testowymi w $H^1$.
Dla stałego $\kappa$ człon ten byłby zerowy mimo
niesymetryczności strumienia. W przykładzie
$\kappa=x_1$ sam operator $A=-\Delta+\partial_2$
jest niesymetryczny: dla $u,v\in C_c^\infty(\Omega)$
różnica $\langle Au,v\rangle-\langle Av,u\rangle$
wynosi $2\int_\Omega(\partial_2u)v\,\mathrm d\lambda_2$.
Wybór $u$ z $\partial_2u\not\equiv0$ i
$v=\partial_2u$ daje ściśle dodatnią różnicę.
Nie istnieje więc różniczkowalny w sensie Frécheta
potencjał o różniczce $Au$: taka różniczka byłaby
ciągła i liniowa względem $u$, a druga różniczka
potencjału byłaby symetryczna.
W przypadku $\kappa(x)=x_1$ minimalizacja samej energii
$\frac12\int_\Omega|\nabla u|^2\,\mathrm d\lambda_2
-\langle f,u\rangle$ nie prowadzi do tego równania.
\end{przyklad}



\begin{uwaga}[Monotoniczność, energia i ograniczenia jednostronne]
\label{rem:pde-v44-minty-obstacle}
Twierdzenie~\ref{thm:pde26-leraylions} nie wymaga istnienia
potencjału: operator monotoniczny może być niesymetryczny,
jak w przykładzie~\ref{ex:pde26-skewflux}. Nie wolno więc
wnioskować, że każde takie równanie jest warunkiem stacjonarności
funkcjonału energii. W rozdziale~\ref{chap:pde29} korzystamy
z innej, dodatkowej struktury: \emph{symetryczna} forma
koercywna określa iloczyn skalarny na $H_0^1(\Omega)$,
a ograniczenie $u\ge\psi$ prowadzi do nierówności wariacyjnej.
Reprezentację jej rozwiązania jako rzutu w normie energetycznej
podamy w stwierdzeniu~\ref{prop:pde-v44-obstacle-projection}.
\end{uwaga}


\section{Porównanie, dodatniość i stabilność względem prawej strony}
\begin{twierdzenie}[Zasada porównania]\label{thm:pde26-comparison}
Niech $f,g\in V^*$, a $u,v\in V$ będą odpowiednio
jedynymi rozwiązaniami $A_pu=f$ i $A_pv=g$.
Jeśli $\langle f-g,\varphi\rangle\le0$ dla wszystkich
nieujemnych $\varphi\in V$, to $u\le v$ prawie wszędzie
w $\Omega$. W szczególności z $f\ge0$ w powyższym
porządku w przestrzeni sprzężonej wynika $u\ge0$ prawie wszędzie.
\end{twierdzenie}
\begin{proof}
Kładziemy $w=(u-v)_+\in V$ i odejmujemy równania słabe:
\[
 \int_{\{u>v\}}(a_p(\nabla u)-a_p(\nabla v))\cdot
 (\nabla u-\nabla v)\,\mathrm d\lambda_d(x)
 =\langle f-g,w\rangle\le0.
\]
Całkowana funkcja jest nieujemna, a przy nierównych
gradientach ściśle dodatnia, na mocy lematu~\ref{lem:pde26-vector}.
Stąd $\nabla w=0$ prawie wszędzie, a nierówność Poincarégo
daje $w=0$. Dla drugiego twierdzenia porównujemy z
zerowym rozwiązaniem równania $A_p0=0$.
\end{proof}

\begin{twierdzenie}[Ilościowa stabilność]\label{thm:pde26-stability}
Niech $u,v$ będą rozwiązaniami z Twierdzenia
~\ref{thm:pde26-plap} dla dowolnych $f,g\in V^*$.
Nie zakładamy uporządkowania $f$ i $g$.
Jeśli $p\ge2$, to
\begin{equation}\label{eq:pde26-holder-data}
 \|u-v\|_V\le
 \bigl((p-1)2^{p-2}\|f-g\|_{V^*}\bigr)^{1/(p-1)}.
\end{equation}
Jeżeli $1<p<2$ i $\max\{\|u\|_V,\|v\|_V\}\le R$,
to istnieje $C_p>0$ takie, że
\begin{equation}\label{eq:pde26-lipschitz-data}
 \|u-v\|_V\le C_p(2R)^{2-p}\|f-g\|_{V^*}.
\end{equation}
W każdym przypadku odwzorowanie $f\mapsto u$
jest ciągłe z $V^*$ do $V$.
\end{twierdzenie}
\begin{proof}
Odejmując równania i testując przez $u-v$ otrzymujemy
\[
 \langle A_pu-A_pv,u-v\rangle
 =\langle f-g,u-v\rangle
 \le\|f-g\|_{V^*}\|u-v\|_V.
\]
Jeżeli $u=v$, oba oszacowania zachodzą natychmiast.
Dalej zakładamy $u\ne v$, więc $\|u-v\|_V>0$.
Dla $p\ge2$ całkujemy~\eqref{eq:pde26-strong},
co daje $2^{2-p}\|u-v\|_V^p/(p-1)$ po lewej stronie,
a zatem~\eqref{eq:pde26-holder-data}.
Dla $1<p<2$ wprowadzamy
$W(x)=|\nabla u(x)|+|\nabla v(x)|$ i
$h=|\nabla u-\nabla v|$. Na zbiorze $W>0$ rozkładamy
$h^p=(W^{p-2}h^2)^{p/2}W^{p(2-p)/2}$
i stosujemy nierówność Höldera z wykładnikami
$2/p$ i $2/(2-p)$:
\[
 \int_\Omega h^p\,\mathrm d\lambda_d
 \le \left(\int_\Omega W^{p-2}h^2\,\mathrm d\lambda_d\right)^{p/2}
 \left(\int_\Omega W^p\,\mathrm d\lambda_d\right)^{(2-p)/2}.
\]
Na zbiorze $W=0$ przyjmujemy $W^{p-2}h^2=0$.
Ponieważ $0\le h\le W$, tak określona funkcja jest
całkowalna, ograniczona przez $W^p$.
Po podniesieniu do potęgi $2/p$ mamy
\[
 \|u-v\|_V^2
 \le \left(\int_\Omega W^{p-2}h^2\,\mathrm d\lambda_d\right)
       \|W\|_p^{2-p}.
\]
Nierówność Minkowskiego daje $\|W\|_p\le2R$.
Z założenia $u\ne v$ wynika $\|W\|_p>0$, więc
$R>0$ i można dzielić przez $\|W\|_p^{2-p}$.
Wykorzystując~\eqref{eq:pde26-subquad}, dostajemy
$\langle A_pu-A_pv,u-v\rangle
\ge(p-1)(2R)^{p-2}\|u-v\|_V^2$.
Dzielimy przez $\|u-v\|_V$ i otrzymujemy
\eqref{eq:pde26-lipschitz-data} z $C_p=1/(p-1)$.
Wreszcie oszacowanie~\eqref{eq:pde26-apriori} zapewnia
wspólne ograniczenie norm rozwiązań, gdy $f,g$
pozostają w ograniczonym zbiorze $V^*$.
\end{proof}

\section{Jawne rozwiązania: kule i pierścienie}
\begin{przyklad}[Rozwiązanie skrętnościowe w kuli]\label{ex:pde26-torsion}
Niech $B_R=\{x\in\R^d:|x|<R\}$, $R>0$, oraz
$\alpha=p/(p-1)$. Definiujemy
\begin{equation}\label{eq:pde26-torsion}
 u(x)=\frac{p-1}{p}\,d^{-1/(p-1)}
 \bigl(R^\alpha-|x|^\alpha\bigr),\qquad x\in\overline B_R.
\end{equation}
Ponieważ $\alpha>1$, funkcja $u$ jest klasy $C^1$
na domkniętej kuli, z gradientem równym zero w środku,
i znika na brzegu. Należy do $W_0^{1,p}(B_R)$:
mnożymy ją przez gładkie odcięcie równe $1$ w odległości
co najmniej $2\varepsilon$ od brzegu i równe $0$
w odległości mniejszej niż $\varepsilon$, o gradiencie
ograniczonym przez $C/\varepsilon$. W warstwie
brzegowej $|u|\le C\varepsilon$, więc błąd w normie
$W^{1,p}$ dąży do zera wraz z miarą warstwy.
Obcięte funkcje przybliżamy splotami o zwartym nośniku.
Dla $x\ne0$ mamy
$\nabla u=-d^{-1/(p-1)}|x|^{\alpha-2}x$, a zatem,
ponieważ $(\alpha-1)(p-1)=1$,
$|\nabla u|^{p-2}\nabla u=-x/d$.
Równość strumieni zachodzi również w zerze przez
ciągłość. Dywergencja pola $-x/d$ wynosi $-1$.
Całkowanie przez części dla funkcji testowych i
gęstość $C_c^\infty(B_R)$ w $W_0^{1,p}(B_R)$
dają $-\Delta_pu=1$ w sensie słabym.
Dla $p>2$ mamy $1<\alpha<2$:
rozwiązanie jest różniczkowalne z ciągłym gradientem,
ale nie należy do $C^2$ w środku kuli:
na dowolnej prostej przez zero, sparametryzowanej
wektorem jednostkowym, drugi symetryczny iloraz
różnicowy w zerze wynosi
$-2d^{-1/(p-1)}|t|^{\alpha-2}/\alpha$,
co dąży do $-\infty$ przy $t\to0$.
To przykład pokazujący, dlaczego nie można automatycznie
przenosić regularności $C^2$ z przypadku liniowego $p=2$.
\end{przyklad}

\begin{przyklad}[Rozwiązanie radialne w pierścieniu]\label{ex:pde26-annulus}
Niech $0<r_0<r_1<\infty$, $d\ge2$ i $u(x)=U(|x|)$
jest radialnym rozwiązaniem słabym $\Delta_pu=0$
w pierścieniu, z $U\in W^{1,p}(r_0,r_1)$.
Współrzędne biegunowe i testy radialne
$v(x)=\psi(|x|)$, $\psi\in C_c^\infty(r_0,r_1)$,
dają po skróceniu przez pole sfery jednostkowej
\[
 \int_{r_0}^{r_1}r^{d-1}|U'|^{p-2}U'\psi'
                 \,\mathrm d\lambda_1(r)=0.
\]
Zatem $r^{d-1}|U'|^{p-2}U'$ jest stałą prawie
wszędzie, gdyż jego pochodna dystrybucyjna znika
na przedziale.
Stąd $U'(r)=c\,r^{-(d-1)/(p-1)}$ z odpowiednio
zdefiniowaną stałą $c\in\R$. Po całkowaniu
\[
 U(r)=\begin{cases}
 a+b\,r^{(p-d)/(p-1)},&p\ne d,\\
 a+b\log r,&p=d.
 \end{cases}
\]
Dla zadanych wartości $U(r_0)=A$, $U(r_1)=B$
i $\gamma=(p-d)/(p-1)\ne0$ otrzymujemy
\[
 U(r)=A+(B-A)\frac{r^\gamma-r_0^\gamma}
                         {r_1^\gamma-r_0^\gamma}.
\]
Dla $p=d$ ułamek zastępujemy przez
$\log(r/r_0)/\log(r_1/r_0)$.
Jest to też granica powyższego ułamka przy
$\gamma\to0$, po podzieleniu licznika i mianownika
przez $\gamma$ i użyciu pochodnej wykładniczej.
Odwrotnie, otrzymane profile mają strumień
$C x/|x|^d$ o zerowej dywergencji w pierścieniu,
więc spełniają równanie dla wszystkich testów,
nie tylko radialnych. Jeżeli $C\ne0$, pole to
jest lokalnie całkowalne również w całej kuli,
lecz jego dywergencja dystrybucyjna wynosi
$C|\mathbb S^{d-1}|\delta_0$: całkowanie przez
części poza kulą promienia $\varepsilon$ daje
wyraz brzegowy $C\int_{\mathbb S^{d-1}}
\varphi(\varepsilon\omega)\,\mathrm d\sigma(\omega)$,
który dąży do $C|\mathbb S^{d-1}|\varphi(0)$.
Dla $C\ne0$ całka z potęgi $p$ gradientu w kuli
promienia $\rho$ jest dodatnią stałą razy
\[
 \int_0^\rho r^{-(d-1)/(p-1)}\,\mathrm d\lambda_1(r).
\]
Jest skończona dokładnie wtedy, gdy $p>d$.
Wtedy profil ma skończoną granicę w zerze, jest
ograniczony i należy do $W^{1,p}$ w kuli:
całkowanie przez części poza kulą promienia
$\varepsilon$ daje dla pierwszych słabych pochodnych
wyraz brzegowy ograniczony przez stałą razy
$\varepsilon^{d-1}$, który znika dla $d\ge2$.
Dla $p\le d$ nie może należeć do $W^{1,p}$,
gdyż jego klasyczny gradient poza zerem nie
jest w $L^p$. Nawet w przypadku $p>d$ profil
nie jest rozwiązaniem jednorodnego równania
w całej kuli, ze względu na obliczony człon Diraca.
\end{przyklad}

\begin{przyklad}[Jednowymiarowe równanie z prawą stroną]\label{ex:pde26-1d}
Niech $\Omega=(0,1)$ i $f\equiv1$. Jednoznaczne rozwiązanie
$-(|u'|^{p-2}u')'=1$ z $u(0)=u(1)=0$ wynosi
\[
 u(x)=\frac{p-1}{p}\left(2^{-p/(p-1)}
 -|x-\tfrac12|^{p/(p-1)}\right).
\]
W istocie strumień $q=|u'|^{p-2}u'$ jest dokładnie
$q(x)=\tfrac12-x$, więc $-q'=1$; warunki brzegowe
wynikają z symetrii profilu. Dla $p=2$ odzyskujemy
$u(x)=x(1-x)/2$. Przynależność $u$ do
$W_0^{1,p}(0,1)$ wynika również z konstrukcji
w Przykładzie~\ref{ex:pde26-torsion} dla $d=1$,
po przesunięciu środka i przyjęciu $R=1/2$.
\end{przyklad}

\begin{center}
\begin{tikzpicture}[x=8.5cm,y=4cm,>=Latex]
 \draw[->] (-.025,0)--(1.05,0) node[right] {$x$};
 \draw[->] (0,-.025)--(0,.66) node[above] {$u(x)$};
 \foreach \x/\lab in {0/0,.5/{1/2},1/1}
  \draw (\x,.009)--(\x,-.009) node[below] {$\lab$};
 \draw[densely dashed,gray] (.5,0)--(.5,.53);
 \draw[blue!75!black,thick,domain=0:1,samples=121]
  plot (\x,{(1/3)*(.5^3-abs(\x-.5)^3)});
 \draw[teal!75!black,thick,domain=0:1,samples=121]
  plot (\x,{.5*(.5^2-abs(\x-.5)^2)});
 \draw[red!75!black,thick,domain=0:1,samples=121]
  plot (\x,{.75*(.5^(4/3)-abs(\x-.5)^(4/3))});
 \draw[blue!75!black,thick] (.7,.58)--(.77,.58)
  node[right,black,font=\small] {$p=3/2$};
 \draw[teal!75!black,thick] (.7,.49)--(.77,.49)
  node[right,black,font=\small] {$p=2$};
 \draw[red!75!black,thick] (.7,.40)--(.77,.40)
  node[right,black,font=\small] {$p=4$};
\end{tikzpicture}

\small Rysunek: profile rozwiązania jednowymiarowego
dla $p=3/2,2,4$ przy prawej stronie równej $1$.
Dla $p>2$ gradient jest ciągły, lecz druga pochodna
w środku przedziału nie istnieje jako skończona liczba.
\end{center}

\section{Pierwsza wartość własna nieliniowego operatora}
\begin{definicja}[Wartość własna $p$-laplasjanu]
\label{def:pde26-eigen}
Liczbę $\lambda\in\R$ nazywamy wartością własną
Dirichleta $p$-laplasjanu na $\Omega$, jeśli istnieje
niezerowa funkcja $u\in W_0^{1,p}(\Omega)$ spełniająca
$-\Delta_pu=\lambda|u|^{p-2}u$ w $V^*$.
Definiujemy też liczbę
\[
 \lambda_{1,p}(\Omega)=
 \inf_{v\in V\setminus\{0\}}
 \frac{\int_\Omega|\nabla v|^p\,\mathrm d\lambda_d(x)}
 {\int_\Omega|v|^p\,\mathrm d\lambda_d(x)}.
\]
\end{definicja}
\begin{twierdzenie}[Osiągalność minimum i równanie własne]
\label{thm:pde26-eigen}
Niech $\Omega\subset\R^d$ będzie ograniczonym obszarem
Lipschitzowskim i $1<p<\infty$.
Wówczas $0<\lambda_{1,p}(\Omega)<\infty$ oraz
istnieje $u\in W_0^{1,p}(\Omega)$ z
$\|u\|_{L^p(\Omega)}=1$, $u\ge0$ prawie wszędzie,
takie, że
\[
 \int_\Omega|\nabla u|^{p-2}\nabla u\cdot\nabla v
 \,\mathrm d\lambda_d(x)=
 \lambda_{1,p}(\Omega)
 \int_\Omega |u|^{p-2}uv\,\mathrm d\lambda_d(x)
 \quad(v\in V).
\]
Dla $r>0$ zachodzi ponadto prawo skalowania
$\lambda_{1,p}(r\Omega)=r^{-p}\lambda_{1,p}(\Omega)$.
\end{twierdzenie}
\begin{proof}
Z nierówności Poincarégo mamy
$\|v\|_p\le C\|\nabla v\|_p$,
więc iloraz Rayleigha jest ograniczony od dołu przez
$C^{-p}>0$. Skończoność otrzymujemy po podstawieniu
dowolnej niezerowej funkcji testowej z $C_c^\infty(\Omega)$.
Bierzemy ciąg minimalizujący z $\|v_n\|_p=1$.
Jego gradienty są ograniczone w $L^p$. Dzięki
refleksywności, po przejściu do podciągu,
$v_n\rightharpoonup u$ w $V$. Zwarte zanurzenie
$W_0^{1,p}(\Omega)\hookrightarrow L^p(\Omega)$
(Twierdzenie~\ref{thm:pde10-Rellich-p})
daje dalszy podciąg silnie zbieżny w $L^p$.
Jego granicą musi być $u$, ponieważ ciągłe
zanurzenie $V\hookrightarrow L^p$ zachowuje
zbieżność słabą. Stąd $\|u\|_p=1$.
Słaba półciągłość normy gradientu dowodzi osiągalności.
Ponieważ $|u|\in V$, $\||u|\|_p=1$ oraz
$\|\nabla|u|\|_p=\|\nabla u\|_p$, możemy wybrać
minimizator nieujemny (wzory na moduł wynikają
z Wniosku~\ref{cor:pde9-truncation}).
Dla ustalonego $v\in V$ kładziemy
$P(t)=\|\nabla(u+tv)\|_p^p$ i $Q(t)=\|u+tv\|_p^p$.
Mamy $Q(0)=1$, zatem $Q(t)>0$ w otoczeniu zera.
Różniczkowanie pod całką uzasadniamy majorantami
z dowodu Twierdzenia~\ref{thm:pde26-minimizer},
odpowiednio dla gradientów i dla samych funkcji.
Otrzymujemy
\[
 \begin{aligned}
 P'(0)&=p\langle A_pu,v\rangle,\\
 Q'(0)&=p\int_\Omega |u|^{p-2}uv\,\mathrm d\lambda_d.
 \end{aligned}
\]
Ponieważ $P/Q$ osiąga minimum w zerze,
$0=P'(0)-\lambda_{1,p}Q'(0)$.
To daje wypisaną tożsamość słabą.
Każda wartość własna $\lambda$ spełnia
$\lambda\ge\lambda_{1,p}$, gdyż testowanie równania
własnego przez niezerową funkcję własną daje
$\lambda=\|\nabla u\|_p^p/\|u\|_p^p$.
Liczba $\lambda_{1,p}$ jest więc rzeczywiście
najmniejszą wartością własną.
Wreszcie podstawienie $v(x)=w(x/r)$
w definicji wartości własnej na $r\Omega$ i
zmiana zmiennych $x=ry$ mnoży licznik przez
$r^{d-p}$, a mianownik przez $r^d$. Skalowanie
jest bijekcją między $W_0^{1,p}(\Omega)$ i
$W_0^{1,p}(r\Omega)$: najpierw sprawdzamy to
dla funkcji testowych, potem stosujemy gęstość
i ciągłość skalowania oraz jego odwrotności. Stąd
prawo skalowania. Nie twierdzimy tutaj, że
nieujemny minimizator jest punktowo ściśle dodatni;
taki wniosek wymaga dodatkowej teorii regularności.
\end{proof}

\section{Monotoniczny schemat różnicowy w jednym wymiarze}
Dzięki jednowymiarowej postaci strumienia możemy
zbudować schemat numeryczny, którego jednoznaczność
wynika z tej samej monotoniczności co twierdzenie
\ref{thm:pde26-plap}. Co istotne, dla stałej prawej
strony nie musimy korzystać z iteracyjnego solvera
nieliniowego.

\begin{twierdzenie}[Dokładna konstrukcja dyskretnego strumienia]
\label{thm:pde26-discrete}
Niech $p>1$ i $N\ge2$ będzie parzystą liczbą całkowitą.
Ustalmy $h=1/N$ oraz $x_i=ih$ dla $0\le i\le N$.
Szukamy liczb $U_0,\ldots,U_N$ spełniających
$U_0=U_N=0$ oraz dla $1\le i\le N-1$ układ
\begin{equation}\label{eq:pde26-discrete}
 -\frac{Q_{i+1/2}-Q_{i-1/2}}h=1,\qquad
 Q_{i+1/2}=a_p\left(\frac{U_{i+1}-U_i}h\right).
\end{equation}
Strumienie są zdefiniowane dla $0\le i\le N-1$,
a w równaniu bilansu $1\le i\le N-1$.
Używamy jednowymiarowej funkcji $a_p$ z Lematu
~\ref{lem:pde26-vector}, w szczególności $a_p(0)=0$.
Układ ma dokładnie jedno rozwiązanie, określone przez
\begin{equation}\label{eq:pde26-flux}
 Q_{i+1/2}=\frac12-\left(i+\frac12\right)h,\qquad
 U_i=h\sum_{j=0}^{i-1}\operatorname{sgn}(Q_{j+1/2})
 |Q_{j+1/2}|^{1/(p-1)}.
\end{equation}
Jeśli $u$ jest rozwiązaniem jawnym z
przykładu~\ref{ex:pde26-1d}, to przy
$\beta=1/(p-1)$ zachodzi
\[
 \max_{0\le i\le N}|U_i-u(x_i)|
 \le C_p h^{\min\{2,1+\beta\}},
\]
gdzie $C_p$ nie zależy od $N$ (dla $p=2$ błąd
w węzłach jest dokładnie zerowy).
\end{twierdzenie}
\begin{proof}
\textbf{Krok 1: bilans strumienia.} Pierwsze równanie
\eqref{eq:pde26-discrete} daje
$Q_{i+1/2}=C-(i+1/2)h$. Odwrotnością funkcji
$a_p(s)=|s|^{p-2}s$ jest
$b_p(q)=\operatorname{sgn}(q)|q|^{1/(p-1)}$.
Zatem $U_{i+1}-U_i=h\,b_p(Q_{i+1/2})$.
Warunek $U_N=U_0$ wymaga
$\sum_{i=0}^{N-1}b_p(C-(i+1/2)h)=0$.
Lewa strona jest ściśle rosnącą i ciągłą funkcją
$C$, a dla $C=1/2$ wyrazy sumy znoszą się parami
przez symetrię względem środka odcinka. Stąd jedyne
$C$ jest równe $1/2$. Formuła~\eqref{eq:pde26-flux}
spełnia więc oba warunki brzegowe i wszystkie równania.

\textbf{Krok 2: niezależny dowód jednoznaczności.}
Dla dwóch rozwiązań $U,V$ pomnóżmy różnicę
równań przez $h(U_i-V_i)$ i zsumujmy po węzłach
wewnętrznych. Oznaczmy $w_i=U_i-V_i$ i $R_i=a_p(D_iU)-a_p(D_iV)$,
gdzie $D_iU=(U_{i+1}-U_i)/h$.
Dyskretne całkowanie przez części ma postać
\[
 -\sum_{i=1}^{N-1}(R_i-R_{i-1})w_i
 =\sum_{i=0}^{N-1}R_i(w_{i+1}-w_i),
\]
co wynika z przesunięcia indeksu i $w_0=w_N=0$.
Stąd
\[
 \sum_{i=0}^{N-1}h\bigl(a_p(D_iU)-a_p(D_iV)\bigr)
 (D_iU-D_iV)=0,
\]
gdzie $D_iU=(U_{i+1}-U_i)/h$. Ścisła monotoniczność
$a_p$ implikuje $D_iU=D_iV$ dla wszystkich $i$,
a zatem $U=V$.

\textbf{Krok 3: zbieżność.} Wzór~\eqref{eq:pde26-flux}
wyraża $U_i$ jako sumę środków przedziałów dla
całki $\int_0^{x_i}b_p(1/2-x)\,\mathrm d\lambda_1(x)$,
która jest równa $u(x_i)$.
Na przedziałach odległych od $1/2$ druga pochodna
funkcji $g(x)=b_p(1/2-x)$ spełnia
$|g''(x)|\le C_p|x-1/2|^{\beta-2}$.
Na przedziale $I$ długości $h$ ze środkiem $m$
rozwinięcie Taylora daje
$|g(x)-g(m)-g'(m)(x-m)|\le
\frac12\sup_I|g''|\,|x-m|^2$.
Po scałkowaniu składnik liniowy znika, a błąd
jest nie większy niż $h^3\sup_I|g''|/24$.
Ponieważ $N$ jest parzyste, środek $1/2$ jest węzłem.
Dla przedziałów o odległości od środka należącej
do $[jh,(j+1)h]$, $j\ge1$, supremum pochodnej
jest nie większe niż $C_p(jh)^{\beta-2}$,
także dla $\beta>2$, bo $(j+1)/j\le2$.
Na każdym poziomie $j$ są co najwyżej dwa przedziały.
Suma błędów jest zatem ograniczona przez
$C_ph^{1+\beta}\sum_{j=1}^{N/2}j^{\beta-2}$.
Dla każdego z dwóch przedziałów przylegających
do środka błąd jest nie większy niż
$h(h/2)^\beta+\int_0^h s^\beta\,\mathrm d\lambda_1(s)$,
czyli co najwyżej $C_ph^{1+\beta}$.
Oszacowanie sumy modułów błędów obowiązuje również
dla dowolnej sumy częściowej, a więc dla każdego węzła.
Dla $\beta<1$ wspomniany szereg jest jednostajnie
ograniczony; dla $\beta>1$ jego suma jest
$O(N^{\beta-1})$, co daje $O(h^2)$.
Gdy $\beta=1$, funkcja $g$ jest afiniczna i metoda
punktu środkowego jest dokładna. Dowód jest zakończony.
\end{proof}

\section{Granice twierdzeń i przypadki niemonotoniczne}
\begin{kontrprzyklad}[Zanik monotoniczności po odjęciu potencjału]
\label{ex:pde26-nonmonotone}
Dla $\lambda>0$ rozważmy
$Bu=A_pu-\lambda|u|^{p-2}u$ w $V^*$,
przy czym drugi składnik rozumiemy przez
$\langle |u|^{p-2}u,v\rangle=
\int_\Omega|u|^{p-2}uv\,\mathrm d\lambda_d(x)$.
Jeśli $\lambda>\lambda_{1,p}(\Omega)$, bierzemy
minimizator $\varphi$ z twierdzenia~\ref{thm:pde26-eigen}.
Wtedy
$\langle B\varphi-B0,\varphi-0\rangle
=\lambda_{1,p}-\lambda<0$, więc $B$ nie jest monotoniczny.
Dla $\lambda=\lambda_{1,p}$ równanie $Bu=0$ ma
co najmniej dwa różne rozwiązania: $u=0$ i $u=\varphi$.
Nie wolno więc automatycznie stosować twierdzenia
Minty'ego--Browdera po dodaniu wyrazu o niekorzystnym znaku.
\end{kontrprzyklad}

\begin{uwaga}[Co zostaje na później?]
Nie udowodniliśmy regularności $C^{1,\alpha}$
ani nierówności Harnacka dla $p$-harmonicznych funkcji;
nie wynikają one z samej monotoniczności. Również
nierówności wariacyjne z przeszkodą wymagają
modyfikacji relacji $Au=f$ na nierówność w zbiorze wypukłym;
wrócimy do nich w rozdziale 29.
Zaproponowana tu metoda Galerkina obejmuje również
strumienie typu Leraya--Lionsa przy właściwych
hipotezach wzrostu, monotoniczności i koercywności,
ale pełnego twierdzenia z zależnością od $x$ i $u$
nie ogłaszamy bez osobnej weryfikacji warunków.
\end{uwaga}


% ROZDZIAŁ 27: reakcja i dyfuzja
